
%=====================================================================
\section{Discussion}
\label{sec:discussion}

\paragraph{What the static network already tells us.}
Conclusions that survive aggregation can be read from \cref{tab:cross}.
For CIT, the final aggregate describes degree heterogeneity, clustering and
connected components. It cannot by itself identify preferential attachment or
ageing: these compare new citations with prior degrees and paper ages.
Even densification requires a sequence of historical $N(t),E(t)$ pairs, rather
than one final graph. For PS and CM, static aggregation preserves the
existence of links but hides their timing, repetition and turnover.

\paragraph{What requires time.}
Some conclusions need windows (Parts I and III) and others the full event order
(Parts II and IV).
For H2, both PS and CM exhibit higher pooled burstiness than their P1
timestamp-shuffled baselines. The excess is relatively small for PS
($+0.0163$) but much larger for CM ($+0.3361$).
However, observed memory is lower than P1 in both networks,
so stronger-than-null positive memory is not supported.
For H3, adjacent-window link overlap also differs from the
randomized baseline, demonstrating that static aggregation
cannot capture temporal link turnover.
Community reorganisation, time-respecting reachability and
spreading speed additionally require temporal information,
although their conclusions depend on the corresponding analyses.

\paragraph{Why network types differ.}
The three networks differ in generating mechanism (physical proximity, voluntary
messaging, citing earlier work), time scale and edge semantics (a contact ends, a
message is a point event, a citation persists).
PS interactions are constrained by school schedules and physical proximity,
whereas CM records voluntary online communication over a much longer period.
These differences may contribute to the stronger excess burstiness in CM.
For H3, PS has lower adjacent-hour Jaccard similarity than P1
($0.2198$ versus $0.3668$), while CM has higher adjacent-week
similarity ($0.0958$ versus $0.0691$).
However, the CM contrast reverses at a 14-day window.
Thus, temporal persistence depends on both network type and
aggregation scale, rather than following a universal pattern.
CIT has persistent edges and measured age-dependent target preferences.
Direct comparisons remain conditional on differences in event definitions
and observation periods.

\paragraph{Robustness.}
We check the choice of window size, the observation period (first vs.\ second
half), the community-detection resolution and the choice of null model.
For H2, PS was evaluated using both raw 20-second contact records
and merged contact episodes. Both definitions produce higher
burstiness and lower memory than P1, although the excess
burstiness changes from $+0.1020$ to $+0.0163$.
For H3, PS shows lower link overlap than P1 across all four
tested windows, while CM changes from higher to lower overlap
as the window increases to 14 days.
The P2b endpoint-swap implementation passes structural invariant
checks, but its mixing behaviour has not been fully assessed.

For CIT, the full-period densification exponent is 1.343 but falls to 1.065 in
the later half, while remaining 1.427 with cutoff 2015. Positive degree preference
persists under all declared variants, although its exponent depends on the degree
offset. The diameter verdict is unchanged by the sampled-source sensitivity
checks. These are qualitative robustness findings, not evidence that coverage
and censoring have been removed.

% Owner: B (measured citation-specific interpretation).
\paragraph{Citation-specific evidence.}
CIT supports degree preference and age-dependent target choice, but separates two
claims often grouped together: cumulative edges grow superlinearly, while the
sampled effective diameter does not shrink. A three-parameter citation curve
fits selected histories well, yet its paired MAE advantage over linear prediction
is unresolved. Annual zero clustering is imposed by the strict-year snapshot
construction. These examples show why representation, fitted-sample description
and held-out prediction require separate interpretations. Late reference coverage,
within-venue selection and right censoring constrain all citation conclusions.

\paragraph{Limitations.}
There is one dataset per network type, so generality is limited and conclusions
hold ``in these data''. For PS and CM, pooled burstiness and memory
may be influenced by heterogeneous edge activity.
P1 preserves the global timestamp distribution and therefore does
not eliminate all temporal nonstationarity.
The short PS observation period also introduces boundary effects
in edge-lifetime measurements.
Window-based Jaccard similarity depends on the chosen aggregation scale.
Although P2b preserves event participation counts, its endpoint swaps
are not guaranteed to sample uniformly from all admissible networks.
The 20-draw null ensembles give coarse tail resolution.

CIT also has within-venue sampling, annual timestamp ties and changing reference
coverage. Its year-block bootstrap intervals are descriptive; cumulative graph
observations are dependent.
Forecast selection excludes low-cited papers, and 38.6\% of short-history WSB fits
hit parameter bounds despite optimizer success in every case; fitted collapse
is not independent validation. The primary age-15 target is fully observed,
whereas 2020 endpoint comparisons mix forecast horizons. Neither removes
within-venue citation-coverage limitations.
Boundary effects and sampling bias remain (finite two-day contact window,
sensor participation).
Citation networks are degenerate temporal networks (DAG, no deletion), which is
why several methods do not apply. Each system also contains several time scales
\citep{holme2015modern}, so a window that suits one process may hide another.
Finally, SocioPatterns data are released under CC BY-NC-SA.
